\documentclass{beamer}
%tlmgr install collection-fontsextra

%%\newlength{\wideitemsep}
%%\setlength{\wideitemsep}{\itemsep}
%%\addtolength{\wideitemsep}{25pt}
%%\let\olditem\item
%%\renewcommand{\item}{\setlength{\itemsep}{\wideitemsep}\olditem}

\usepackage{caption}
\usepackage{hyperref}
\usepackage{verbatim}  % allows begin...end{comment}
\usepackage{graphicx}
\usepackage[skins,raster]{tcolorbox}
\usepackage{bbm}
%\usepackage{natbib}
\usepackage{hyperref}
\hypersetup{colorlinks=true, linkcolor=black, urlcolor=blue, citecolor=red}


%%%%%%%%%%%%%%%%%%%
%% Beamer Set-up
%%%%%%%%%%%%%%%%%%%
\newcommand{\footer}{Dissertation Work Review}

\mode<presentation>{\usetheme{CambridgeUS}} % many others are available
\setbeamercolor{title}{bg=crimson,fg=white}
\setbeamertemplate{page number in head/foot}{} % Include/Remove page numbering
\setbeamerfont{footnote}{size=\tiny}
\setbeamertemplate{navigation symbols}{} % eliminate nav symbols in footer


\definecolor{crimson}{HTML}{A51C30}
\definecolor{darkred}{HTML}{A51C30}
\definecolor{offwhite}{HTML}{A1A1A1}

\newcommand{\question}[1]{\textcolor{crimson}{#1}}

\setbeamertemplate{footline}
{
  \leavevmode%
  \hbox{%
  \begin{beamercolorbox}[wd=.333333\paperwidth,ht=2.25ex,dp=1ex,center]{author in head/foot}
    \usebeamerfont{author in head/foot}\insertshortauthor%~~\beamer@ifempty{\insertshortinstitute}{}{(\insertshortinstitute)}
  \end{beamercolorbox}%
  \begin{beamercolorbox}[wd=.333333\paperwidth,ht=2.25ex,dp=1ex,center]{title in head/foot}%
    \usebeamerfont{title in head/foot}\footer 
  \end{beamercolorbox}%
  \begin{beamercolorbox}[wd=.333333\paperwidth,ht=2.25ex,dp=1ex,right]{date in head/foot}%
    \usebeamerfont{date in head/foot}\insertshortdate{}\hspace*{2em}
    \insertframenumber{} / \inserttotalframenumber\hspace*{2ex} 
  \end{beamercolorbox}}%
  \vskip0pt%
}

\useoutertheme{miniframes}

\AtBeginSection[]
{
  \begin{frame}<beamer>
%    \frametitle{Outline for Section \thesection}
    \frametitle{Outline}
    \tableofcontents[currentsection]
  \end{frame}
}


%%%%%%%%%%%%%%%%%%%%%%%%
% Bibliography Options 
%%%%%%%%%%%%%%%%%%%%%%%%

\usepackage[backend=biber, style=alphabetic, url=false, doi=false, maxcitenames=1, isbn=false]{biblatex}
%\addbibresource{ReadingGroup.bib}
%\usepackage[backend=biber, giveninits=true]{biblatex}


\title{Dissertation Work Formulation}
\author{Dominic DiSanto; Tianxi Cai \& Junwei Lu}
\institute[]{Department of Biostatistics, Harvard University}
\date{\today}


% Custom commands 
\newcommand{\Loss}{\mathcal{L}}
\newcommand{\R}{\mathbb{R}}
\newcommand{\EE}{\mathbb{E}}
\newcommand{\Exp}{\EE}
\newcommand{\Xbf}{\mathbf{X}}
\newcommand{\Wbf}{\mathbf{W}}
\renewcommand{\Pr}{\mathbb{P}}
\newcommand{\ATE}{\tau}
\newcommand{\tATE}{\text{ATE}}

\newcommand{\Xb}{\mathbf{X}}
\newcommand{\Yb}{\mathbf{Y}}
\newcommand{\Tb}{\mathbf{T}}

\newcommand{\Zbar}{\overline{Z}}
\newcommand{\Abar}{\overline{A}}
\newcommand{\Lbar}{\overline{L}}

\newcommand{\zbar}{\overline{z}}
\newcommand{\abar}{\overline{a}}
\newcommand{\lbar}{\overline{\ell}}


\newcommand{\IPW}{{\text{IPW}}}
\newcommand{\AIPW}{{\text{AIPW}}}
\newcommand{\DIPW}{{\text{DIPW}}}
\newcommand{\ORA}{{\text{ORA}}}
\newcommand{\Data}{\mathcal{D}}

\newcommand{\subitem}[1]{\begin{itemize} \item #1 \end{itemize}}

\newcommand{\widesubitem}[2]{\begin{itemize} \item #1 \end{itemize} \vspace{#2}}

\newcommand{\PP}{\mathbb{P}}
\newcommand{\RR}{\mathbb{R}}
\newcommand{\NN}{\mathbb{N}}

\newcommand{\Ind}{\mathbbm{1}}

\newcommand{\pistar}{\pi^*}

\newcommand{\pconv}{\stackrel{p}{\rightarrow}}
\newcommand{\dconv}{\stackrel{d}{\rightarrow}}

\newcommand{\assign}{\mathbbm{d}}

\renewcommand{\texttilde}{\textasciitilde}
\newcommand\indp{\protect\mathpalette{\protect\indpT}{\perp}}
\def\indpT#1#2{\mathrel{\rlap{$#1#2$}\mkern2mu{#1#2}}}


\newenvironment{wideitemize}{\itemize\addtolength{\itemsep}{8pt}}{\enditemize}





%%%%%%%%%%%%%%%%%%%%%%%%
    % START %%%%%%%%%%%%
%%%%%%%%%%%%%%%%%%%%%%%%
\begin{document}

\frame[plain]{\titlepage}

\begin{frame}{Today's Goals/Questions}
    \begin{itemize}\setlength\itemsep{0.75em}
        \item Review our set-up (observables, data-generating processes)
            \subitem{Longitudinal modified treatment policy (LMTP) framework}
        \item Understand most relevant extensions of treatment assignment
        \item Understand PheWAS-relevant formulation of outcomes in $Y$
    \end{itemize}
\end{frame}


\section{Set-up}


\begin{frame}{Set-Up and Observables}
    \begin{itemize}\setlength\itemsep{1em}
        \item For times $t=1, \dots, \tau$, observe $Z = (L_1, A_1, \cdots L_\tau, A_\tau, Y) \sim \PP$ \vspace{0.5em}
            \subitem{Permit the short-hand $\Abar_t = (A_t, \cdots, A_1)$, $\Lbar_t$ similarly defined, and history as $H_t = (\Abar_{t-1}, \Lbar_t)$} \vspace{0.5em}
        
            \subitem{Covariates $L_t, \ e.g. 
            \begin{cases} 
                \text{EHR Codes Accumulation or general Hawkes process} \\ 
                \text{Demographic/static information} \\
            \end{cases}
            $
            } \vspace{0.5em}

            \subitem{Treatment $A_t$: binary, multinomial, continuous (e.g. dose)
                \subitem{More detailed discussion under Treatment/Policy Evaluation section}
            } \vspace{0.5em}

            \subitem{Outcome $Y_t$
                \subitem{For first discussion, assume univariate $Y$} 
                \subitem{More detailed discussion under PheWAS section}
            } 
        \end{itemize}

\end{frame}

\begin{frame}{{Data Generating Process \& Treatment Mechanism}}

    \begin{itemize}\setlength\itemsep{1em}
        \item Assume generating functions with noise $U = (U_A, U_L, U_Y)$

            \widesubitem{$A_t = f_{A_t}(H_t,  U_{A_t})$}{0.5em}
            \widesubitem{$L_t = f_{L_t}(A_t, H_t,  U_{L_t})$}{0.5em}
            \widesubitem{$Y = f_Y(A_{\tau}, H_\tau,  U_Y)$}{0.5em}

        \item Assume a new assignment function $\assign(a_t, h_t^\assign) \rightarrow a^\assign_t $

        \item Typical non-parametric SEM set-up (similar, somewhat more general to Aaron Sonabend's SSRL set-up)
    \end{itemize}
\end{frame}

\begin{frame}[allowframebreaks]{Example $\assign$ functions}

    \begin{itemize}\setlength\itemsep{1.5em}
        \item Thresholding, e.g. patient's with SBP$>130$ receive drug $H$ and otherwise drug $G$ (here $SBP = L_t$)
        \subitem{$A_t^\assign = \assign(a_t, h_t) = \Ind(SBP_t>130)H + \Ind(SBP_t\leq 130)G$}

        \item Shifting for continuous treatments, e.g. drug dose 
        \subitem{If $P(A_t < u_t(h_t) | H_t = h_t) = 1$, then 
        \begin{gather*}
            \assign(a_t, h_t) 
            = 
            \begin{cases}
                a_t + \delta &  a_t \leqslant u_t(h_t) - \delta \\ 
                a_t & a_t > u_t(h_t)-\delta
            \end{cases}
        \end{gather*}
        } 

        \newpage 

        \item Can induce stochastic interventions using noise $\varepsilon\indp\PP$, now $\assign(a_t, h_t, \varepsilon)$
        \item 
        \subitem{Ex: Shifted/Incremental Propensity Score. For a binary treatment with density $g(a|h_t)$ (and user defined $\delta>0$),}

        \begin{gather*}
            g_t^\assign(1|h_t) = 
            \frac{\delta g_t(1|h_t)}{\delta g_t(1|h_t) + 1-g_t(1|h_t)}
            \\ \\ 
            \assign(a_t, h_t) = \Ind(\epsilon_t \leq g_t^\assign(1|h_t))
            \text{ for } \epsilon_t \sim U(0,1)
        \end{gather*}



    \end{itemize}


\end{frame}

\begin{frame}{Counterfactuals}
    \begin{itemize}

    \item Observe $\Abar_t$, with counterfactuals generated by the assignment functions (defined recursively through $t=1$)

        \widesubitem{Write the counterfactual $A_t(A^\assign_{t-1}) = f_{A,t}(H_t(\Abar_{t-1}^\assign), U_{A,t})$
            \subitem{``What treatment would we observe at $A_t$, had we applied our intervention policy $\assign$ through time $t-1$?''}
            \subitem{``Natural value of treatment''}
        }{0.3em}

        \widesubitem{$L_t(\Abar^\assign_{t-1}) = f_{L,t}(A_{t-1}^\assign, H_{t-1}^\assign, U_{L,t})$}{0.3em}

        \widesubitem{$Y(\Abar^\assign) = f_Y(\Abar_\tau^\assign, H_\tau(\Abar_{\tau-1}^\assign), U_Y)$}{0.3em}

    \item Estimand under a given policy $\assign$ is $\theta := \EE\left[ Y(\Abar^\assign) \right]$

    \end{itemize}
\end{frame}



\begin{frame}{Extensions/Problem Formulation}
    \begin{itemize}\setlength\itemsep{1.5em}
        \item Extend the framework's treatment regime (for a simple/univariate outcome)
        \widesubitem{Via assignment function(s) $\assign$}{1em}
        \widesubitem{Via treatment vector $A_t$}{1em}
        %\item Nima has a ``causal survival'' paper extending LMPT that includes competing risks 

        \item For simple treatment/policy comparisons, multivariate or high-dimensional $Y\in\RR^d$ 
    \end{itemize}
\end{frame}

\section{Treatment/Policy Evaluation ($\assign, A$)}



\begin{frame}{$\assign$}
    \begin{itemize}\setlength\itemsep{1em}
        \item We can alter $\assign$ 
        \subitem{Conceptually we are then trying to estimate parameters/functions that describe treatment patterns, while comparing $\EE[Y^\assign]$}
        \subitem{We can write assignment function as complicated/high-dimensional function 
        \widesubitem{$\assign(\abar_t, \lbar_t; \alpha, \beta) = \alpha^T\abar_t + \beta^T\lbar_t$}{0.5em}
        \widesubitem{Non-parametric $\assign = m(\abar_t, \lbar_t)$}{0.5em}
        }

        \item We can extend $A_t$ (e.g. $A_t$ is now multivariate, captures treatment more globally)
            \widesubitem{Conceptually we still have ``control'' over assignment via $\assign$ that we can (and must) specify}{0.5em}
            \widesubitem{e.g. high-dimensional, categorical $A\in\RR^d$ drug choices}{0.5em}

    \end{itemize}
\end{frame}



\section{PheWAS ($Y$)}

\begin{frame}{PheWAS}
    \begin{itemize}\setlength\itemsep{1em}
        \item Can simplify treatment w/in previous framework
        \subitem{i.e. binary treatment $A$, two policies $\assign_0(a_t) = a_t$, $\assign_1(a_t) = 1-a_t, t=\{0, 1\}$}
        \item Now interested in settings with outcome diversity, $Y \in \RR^d$
        \item Sub-group analysis? 
            \subitem{e.g. In RA, how could we identify sub-types of heart failure (preserved vs reduced ejection fraction) for which anti-TNF is effective?}

        \item Disjoint outcomes? 
        \subitem{Across $Y\in\RR^d-$outcomes (possibly dependent, overlapping), how do we identify the subset of $Y$ (if any) for which drug $G$ has non-0 treatment effect?}
    \end{itemize}
\end{frame}


\section{Summary/Questions}

\begin{frame}{Summary/Questions}
    \begin{itemize}\setlength\itemsep{1.5em}
        \item Is this (LMTP) set-up the appropriate framework which we can extend? 
        \item What are the relevant extensions of treatment trajectories via $A_t, \assign$?
        \subitem{Parameterizing $\assign(a_t, h_t)$ as a high-dimensional policy}
        \widesubitem{Specifying a form for $A_t$ with simple policy $\assign$}{0.5em}

        \item What are the relevant extensions outcomes $Y$ for desired PheWAS studies (and how can we do better than na\"ive multiple comparisons adjustment)?
        \subitem{Sub-group identification?}
        \widesubitem{High-/Multi-dimensional $Y\in\RR^d$?}{0.5em}

    \end{itemize}
\end{frame}


\title{Appendix}
\author{} 
\institute{} 
\subtitle{} 
\date{} 
\frame{\titlepage}

\appendix 
    \section{Identification Assumptions} 


    \begin{frame}{Identification}
        \begin{enumerate}
            \item Positivity: $(a_t, h_t) \in \text{Supp}(A_t, h_t) \Rightarrow \assign(a_t, h_t) \in \text{Supp}, \forall t \in [\tau]$ 

            \item Strong Sequential Randomization: $U_{A_t} \indp \underline{U}_{L, t+1}, \underline{U}_{A, t+1} \mid H_t, \ \forall t \in[\tau]$
        \end{enumerate}
        \vspace{2em}

        The Strong Sequential Randomization assumption can be slightly weakened for stochastic interventions

        %Estimation in general setting proceeds via TMLE (with $\tau+1$ robustness)
    \end{frame}

    
    \iffalse 
    \begin{frame}{Estimation under LMTP (Simple Set-up)}
    \end{frame}
    \fi 

    \section{Misc. Goals/Timeline}

    \begin{frame}[allowframebreaks]
    \frametitle{\large{Rough Timeline of G3 ('24-Spring '25)}}

        \begin{itemize}\setlength\itemsep{1em}
            \item Continue/complete early work 
            \begin{itemize}
            \item \textbf{August} Submit Weijing-Zongqi applied paper 
            \item \textbf{Fall-Winter} Complete draft of optimal assortment work with Junwei
            \end{itemize}
            \item \textbf{October -} Finalize committee 
            \begin{itemize}
                \item Possibly Rajarshi, Sebastien, Nima 
            \end{itemize}
            \item \textbf{March -}  Oral Examination
        \end{itemize}
    \end{frame}



\end{document}0